
\section{Grundlagen und gängige C2-Frameworks}
\subsection{Command-and-Control (C2) - Grundkonzept}
\subsubsection{Definition und Zweck von C2}
Command-and-Control (C2) bezeichnet im Kontext von Cyberangriffen die Gesamtheit der Mechanismen und Infrastrukturen, 
die es einem Angreifer ermöglichen, kompromittierte Systeme zentral zu steuern. 
Nach erfolgreicher Kompromittierung eines Zielsystems dient eine C2-Infrastruktur dazu, Befehle an Implantate oder Schadsoftware zu übermitteln, 
Rückmeldungen zu empfangen und weitere Aktionen zu koordinieren. 
C2 stellt somit die kommunikative Verbindung zwischen Angreifer und Zielsystem dar und ist eine wesentliche Voraussetzung für nachhaltige und zielgerichtete Angriffe.\cite{c2_varonis2026}
\subsubsection{Beaconing und Kommunikation in C2-Frameworks}
Die Kommunikation zwischen kompromittierten Systemen und einer Command-and-Control-Infrastruktur bildet einen zentralen Bestandteil moderner C2-\ Frameworks. 
Ziel ist es, Befehle und Rückmeldungen zuverlässig sowie möglichst unauffällig zu übertragen. 
Zu diesem Zweck unterstützen aktuelle Frameworks unterschiedliche Kommunikationsprotokolle, die sich an legitimen Netzwerkverkehr anlehnen und verschiedene Einsatzszenarien abdecken.
\newline\newline
Ein verbreitetes Kommunikationsmodell ist das Beaconing, bei dem kompromittierte Systeme in regelmäßigen oder adaptiven Intervallen selbstständig eine Verbindung zum C2-Server herstellen. 
Dieses Modell begünstigt ausgehenden Netzwerkverkehr und ist insbesondere in Umgebungen mit restriktiven Firewall-Regeln von Vorteil.
Demgegenüber stehen sessionbasierte Modelle mit persistenten Verbindungen, die eine unmittelbare Interaktion ermöglichen, jedoch eine erhöhte Erkennbarkeit aufweisen können. \cite{beacon_session}
\newline\newline
Zur Umsetzung dieser Modelle nutzen C2-Frameworks verschiedene Protokolle. 
Häufig kommt TLS-geschützter HTTP(S)-Verkehr zum Einsatz, da dieser in den meisten Netzwerken erlaubt ist und sich gut in regulären Web-Traffic integrieren lässt. 
Ergänzend werden DNS-basierte Mechanismen verwendet, um auch in stark eingeschränkten Netzwerken eine grundlegende Kommunikation zu ermöglichen. 
Moderne Frameworks unterstützen zudem verschlüsselte Tunnelverfahren wie WireGuard, die eine robuste und performante Kommunikation bei gleichzeitig hohem Sicherheitsniveau bieten.
\newline\newline
Die detaillierte Betrachtung dieser Kommunikationsmechanismen sowie deren operative Auswirkungen erfolgt in den nachfolgenden Abschnitten dieser Arbeit.
\subsection{Anforderungen an moderne C2-Frameworks}
Moderne Command-and-Control-Frameworks müssen eine Vielzahl technischer und operationeller Anforderungen erfüllen, um in realistischen Angriffsszenarien einsetzbar zu sein. 
Insbesondere im Kontext professioneller Sicherheitsprüfungen ist es erforderlich, dass C2-Infrastrukturen sowohl funktional leistungsfähig als auch anpassbar und kontrollierbar sind. 
Die folgenden Kriterien stellen zentrale Anforderungen dar, die bei der Bewertung moderner Frameworks berücksichtigt werden müssen.
\subsubsection{Stealth und OPSEC}
Eine der wesentlichsten Anforderungen an moderne C2-Frameworks ist ein hohes Maß an Stealth. Darunter ist die Fähigkeit zu verstehen, 
Kommunikations- und Laufzeitverhalten so zu gestalten, dass sie möglichst wenig Auffälligkeiten gegenüber Sicherheitsmechanismen erzeugen. 
In modernen IT- und insbesondere OT-Umgebungen werden Netzwerkverkehr, Prozessverhalten und Speicheraktivitäten zunehmend überwacht. 
Auffällige Muster können daher schnell zu einer Detektion führen.
\newline\newline
Um diesem Risiko zu begegnen, müssen C2-Frameworks in der Lage sein, ihre Kommunikation und ihr Verhalten an die jeweilige Zielumgebung anzupassen. 
Dies umfasst die Nutzung von legitimen Protokollen, die Integration in regulären Netzwerkverkehr sowie die Vermeidung von auffälligen Mustern in der Prozess- und Speicheraktivität. 
Darüber hinaus ist eine sorgfältige OPSEC (Operational Security) unerlässlich, um die Angriffsinfrastruktur vor Entdeckung zu schützen und die Effektivität der Simulation zu gewährleisten.
\subsubsection{Modularität}
Modularität ist ein weiteres zentrales Kriterium für moderne C2-Frameworks. 
Ein modularer Aufbau ermöglicht es, verschiedene Komponenten wie Kommunikationsprotokolle, Verschlüsselungsmechanismen oder Implantattypen flexibel zu kombinieren und an spezifische Anforderungen anzupassen. 
Dies erhöht nicht nur die Einsatzmöglichkeiten, sondern erleichtert auch die Wartung und Weiterentwicklung der C2-Infrastruktur.
\subsubsection{Umgehung von EDR/AV-Systemen}
Die zunehmende Verbreitung von Endpoint Detection and Response (EDR) sowie verhaltensbasierten Antivirenlösungen verschärft die Anforderungen an C2-Frameworks zusätzlich.
Moderne Sicherheitsprodukte analysieren Prozesse, Speicherzugriffe und Netzwerkaktivitäten anhand heuristischer Modelle und Anomalieerkennung\cite{heuristic}. Vor diesem Hintergrund genügt es nicht, lediglich bekannte Signaturen zu vermeiden.
\newline\newline
Vielmehr muss ein C2-Framework in der Lage sein, generische Erkennungsmerkmale zu minimieren. 
Dies betrifft sowohl die Art der Prozessausführung als auch das Kommunikationsverhalten. 
Die zuvor beschriebene Modularität unterstützt dies, da sie es erlaubt, Mechanismen situationsabhängig anzupassen. 
Stealth, OPSEC und Modularität wirken somit zusammen, um die Wahrscheinlichkeit einer Detektion zu reduzieren und gleichzeitig realistische Angriffsszenarien zu ermöglichen.
\subsubsection{Flexibilität und Anpassungsfähigkeit}
Die genannten Anforderungen münden schließlich in der übergeordneten Notwendigkeit der Flexibilität. IT- und insbesondere OT-Umgebungen unterscheiden sich stark hinsichtlich Netzwerksegmentierung, Sicherheitsniveau und regulatorischer Rahmenbedingungen.
Ein C2-Framework muss daher in der Lage sein, sich an unterschiedliche technische und organisatorische Gegebenheiten anzupassen.
\newline\newline
Flexibilität umfasst sowohl die Unterstützung verschiedener Betriebssysteme und Kommunikationsprotokolle als auch die Möglichkeit, unterschiedliche Angriffsszenarien abzubilden.
Während in einer Umgebung kurzfristige Interaktion im Vordergrund stehen kann, ist in einer anderen die Simulation persistenter, langfristiger Zugriffe entscheidend. 
Die Fähigkeit, Kommunikationsmodelle, Module und Einsatzparameter situationsabhängig zu konfigurieren, ist daher essenziell.
\newline\newline
Gerade im regulierten Umfeld, etwa im Kontext von NIS-2 oder kritischen Infrastrukturen, ist eine kontrollierte und nachvollziehbare Anpassungsfähigkeit von Bedeutung. 
Ein modernes C2-Framework muss nicht nur technisch leistungsfähig sein, sondern auch strukturiert einsetzbar, dokumentierbar und in bestehende Sicherheitsprozesse integrierbar bleiben.
\subsection{Überblick über gängige C2-Frameworks}
\subsubsection{Metasploit (MSF)}
Metasploit (MSF) zählt zu den bekanntesten und am weitesten verbreiteten Frameworks in der Sicherheitscommunity. Ursprünglich als Plattform zur Entwicklung und Integration von Exploits konzipiert, hat es sich zu einem umfassenden Werkzeug für Penetrationstests und angriffsnahe Sicherheitsanalysen entwickelt. Das Framework stellt eine Vielzahl von Modulen bereit, darunter Exploits, Payloads, Encoders und Auxiliary-Komponenten, und deckt damit insbesondere die Phase der Initialkompromittierung strukturiert ab \cite{msf_module}. 
Command-and-Control-Funktionalitäten werden primär über das Meterpreter-Modul realisiert.
\newline\newline
Meterpreter ermöglicht eine interaktive, sessionbasierte Steuerung kompromittierter Systeme und unterstützt unter anderem Dateioperationen, Prozessinteraktionen sowie Speicherzugriffe. Die unmittelbare und latenzarme Ausführung von Befehlen ist insbesondere in klassischen Penetrationstests vorteilhaft, da Schwachstellen schnell validiert und Auswirkungen anschaulich demonstriert werden können. Gleichzeitig ist zu berücksichtigen, dass persistente Sessions im Vergleich zu Beaconing-basierten Modellen eine erhöhte Netzwerkauffälligkeit verursachen können. 
Dauerhafte Verbindungen oder interaktive Shells lassen sich in überwachten Umgebungen leichter identifizieren, insbesondere bei Einsatz moderner EDR- oder Netzwerküberwachungssysteme. 
Die hohe Interaktivität stellt somit einen operativen Vorteil dar, kann jedoch in realitätsnahen Red-Team-Szenarien mit erhöhten Detektionsrisiken verbunden sein.
\newline\newline
Im Unterschied zu speziell für langfristige Red-Team-Operationen entwickelten Frameworks liegt der Schwerpunkt von Metasploit weniger auf ausgeprägter OPSEC oder adaptiven Kommunikationsmodellen. Standardisierte Payloads und bekannte Artefakte führen in produktiven Umgebungen nicht selten zu einer schnellen Erkennung. Zwar bestehen Anpassungs- und Erweiterungsmöglichkeiten, jedoch war eine langfristig unauffällige C2-Infrastruktur nicht primärer Designfokus des Frameworks.
\newline\newline
Aus wirtschaftlicher Sicht ist zwischen dem frei verfügbaren Metasploit Framework und kommerziellen Varianten wie Metasploit Pro zu unterscheiden. Während das Open-Source-Framework kostenfrei genutzt werden kann, bieten die kostenpflichtigen Versionen zusätzliche Funktionen, etwa im Bereich automatisierter Berichterstellung, Web-basierter Verwaltung und Projektmanagement. Diese werden üblicherweise im Rahmen jährlicher Lizenzmodelle angeboten und können regionalen Prüfungen unterliegen. 
Für Organisationen ergibt sich daher eine Abwägung zwischen funktionalem Mehrwert und Lizenzkosten.
\subsubsection{Cobalt Strike}
Cobalt Strike wurde gezielt für professionelle Red-Team-Operationen entwickelt und hat sich in den vergangenen Jahren als etablierter Standard im kommerziellen Red-Team-Umfeld positioniert. 
Im Gegensatz zu Metasploit, dessen Ursprung in der strukturierten Exploit-Integration liegt, steht bei Cobalt Strike die durchführung fortgeschrittener Angriffe über mehrere Phasen der Kill Chain hinweg im Vordergrund. 
Der Fokus liegt damit weniger auf der initialen Schwachstellenausnutzung als auf nachhaltiger Post-Exploitation und kontrollierter Steuerung kompromittierter Systeme.
\newline\newline
Zentrales Element der Architektur ist das sogenannte Beacon-Konzept. Anders als primär sessionbasierte Modelle setzt Cobalt Strike auf eine adaptive Beaconing-Kommunikation, die über verschiedene Transportmechanismen wie HTTP(S), DNS oder SMB erfolgen kann. Durch konfigurierbare Intervalle, verschleierte Datenstrukturen und flexible Infrastrukturparameter lässt sich das Kommunikationsverhalten an die jeweilige Zielumgebung anpassen. 
Dieses Modell reduziert im Vergleich zu persistenten Sessions die Netzwerkauffälligkeit und eignet sich insbesondere für längerfristige, realitätsnahe Angriffssimulationen.
\newline\newline
Neben der Kommunikationsarchitektur bietet Cobalt Strike umfangreiche Funktionen zur Teamkoordination. Mehrere Operatoren können parallel auf einem Teamserver arbeiten, wodurch komplexe Angriffsszenarien strukturiert geplant und durchgeführt werden können. Ergänzt wird dies durch integrierte Mechanismen zur Verwaltung von Beacons, zur Ausführung von Post-Exploitation-Aktivitäten sowie zur Anpassung operativer Parameter. 
Diese Kombination aus technischer Flexibilität und organisatorischer Struktur macht das Framework besonders leistungsfähig im professionellen Einsatz.
\newline\newline
Gleichzeitig ist zu berücksichtigen, dass Cobalt Strike aufgrund seiner weiten Verbreitung verstärkt im Fokus von Sicherheitsanbietern steht. Standardisierte Konfigurationen oder unzureichend angepasste Profile können daher ebenfalls zu einer schnellen Detektion führen. Außerdem ist die kommerzielle Lizenzierung mit jährlichen Kosten verbunden, die je nach Umfang der Nutzung und regionalen Gegebenheiten variieren können. 
Für Organisationen, die auf professionelle Red-Team-Operationen setzen, stellt sich daher die Frage nach dem Kosten-Nutzen-Verhältnis im Vergleich zu alternativen Lösungen.
\subsubsection{Sliver}
Sliver stellt eine moderne, Open-Source-basierte Command-and-Control-Lösung dar, die von Beginn an mit dem Ziel entwickelt wurde, flexible und modulare Red-Team-Infrastrukturen bereitzustellen. Im Unterschied zu Metasploit, das seinen Schwerpunkt auf Exploit-Integration legt, und Cobalt Strike, das als kommerzielles Produkt stark auf koordinierte Red-Team-Operationen ausgerichtet ist, positioniert sich Sliver als technisch leistungsfähige, zugleich kostenfreie Alternative mit infrastrukturellem Fokus.
\newline\newline
Architektonisch unterstützt Sliver sowohl Beaconing- als auch sessionbasierte Kommunikationsmodelle und bietet eine breite Auswahl an Transportprotokollen, darunter TLS-gesicherter HTTP(S)-Verkehr, DNS sowie WireGuard-basierte Tunnel. Diese Vielfalt ermöglicht eine flexible Anpassung an unterschiedliche Netzwerkbedingungen und Sicherheitsanforderungen. Besonders hervorzuheben ist die plattformübergreifende Ausrichtung des Frameworks, das Implantate für verschiedene Betriebssysteme generieren kann und damit heterogene Zielumgebungen adressiert.
\newline\newline
Ein wesentliches Merkmal von Sliver ist die konsequente Modularität der Architektur. Kommunikationsmechanismen, Implantate und Erweiterungen sind klar strukturiert voneinander getrennt, was eine gezielte Anpassung einzelner Komponenten erlaubt. Darüber hinaus unterstützt Sliver einen Multi-Operator-Betrieb, wodurch parallele Arbeit mehrerer Tester innerhalb einer gemeinsamen Infrastruktur möglich ist. Diese Eigenschaften rücken das Framework funktional näher an kommerzielle Lösungen heran, ohne an ein Lizenzmodell gebunden zu sein.
\newline\newline
Aus wirtschaftlicher Sicht stellt der Open-Source-Charakter einen signifikanten Unterschied zu Cobalt Strike dar. Es fallen keine Lizenzkosten pro Operator oder Server an, was insbesondere für interne Sicherheitsabteilungen oder Organisationen mit mehreren Testumgebungen von Vorteil sein kann. Gleichzeitig liegt die Verantwortung für Wartung, Konfiguration und Infrastrukturdesign vollständig beim betreibenden Team, da kein kommerzieller Support im klassischen Sinne vorgesehen ist.
\newline\newline
Im Vergleich zu Metasploit ist Sliver weniger auf die initiale Ausnutzung von Schwachstellen fokussiert, sondern stärker auf die strukturierte Steuerung und Verwaltung kompromittierter Systeme ausgelegt. Im Vergleich zu Cobalt Strike bietet es ähnliche konzeptionelle Ansätze hinsichtlich Beaconing und Teamfähigkeit, jedoch in einem offenen und anpassbaren Entwicklungsmodell.
\newline\newline
Die detaillierte Analyse der Architektur, Kommunikationsmechanismen sowie der operativen Eigenschaften von Sliver erfolgt im nachfolgenden Kapitel. Dort wird insbesondere untersucht, inwiefern das Framework die zuvor definierten Anforderungen an moderne C2-Infrastrukturen erfüllt.